\documentclass{article}
\usepackage{graphicx} % Required for inserting images
\usepackage{amsmath}
\usepackage{amssymb}
\usepackage{amsfonts}
\usepackage{verbatim}
\usepackage{stmaryrd}
%\usepackage{bussproofs}
\usepackage{mathabx}
\usepackage{color}
\usepackage{tikz}
\usetikzlibrary {automata,positioning}

\newcommand{\setenc}[1]{\llparenthesis #1 \rrparenthesis}
\newcommand{\set}[1]{\{#1\}}
\newcommand{\tuple}[1]{\langle #1\rangle}
\newcommand{\true}{\texttt{true}}
\newcommand{\false}{\texttt{false}}
\renewcommand{\diamond}{\lozenge} %%\diamondsuit
\newcommand{\Next}{\ovoid} %{\circ}
\newcommand{\Nat}{\mathbb{N}}
\newcommand{\sem}[1]{\llbracket #1\rrbracket}
\newtheorem{property}{Property.}

\title{COSC437 Security Protocols: Assignment 2}
\author{Name: Navaneethakrishna Sridhar \\ Student ID: 12590145}
\date{1 Sep 2026}

\begin{document}

\maketitle

\paragraph{The model.}
I put the model in \texttt{bwbakery.tla}. It is Algorithm~2 from~\cite{taubenfeld}, for three processes, in PlusCal, then translated into the Toolbox. Same shape as the week~5 Peterson demo: \texttt{process (P \textbackslash in Procs)} inside \texttt{ncs: while (TRUE)}, with \texttt{N <- 3} and \texttt{Procs == 1..N} for the paper's $1..n$.

Shared variables are the ones from the paper: \texttt{color}, \texttt{choosing}, and the two halves of a coloured ticket, \texttt{mycolor} and \texttt{number}. I split those two because line~2 writes the colour and line~3 writes the number. Labels follow the paper: \texttt{e1}--\texttt{e4} doorway (lines~1--4), \texttt{e5}--\texttt{e10} waiting loop (lines~5--10), \texttt{cs} line~11, \texttt{x12}/\texttt{x13} exit. \texttt{choosing} and \texttt{number} start at \texttt{FALSE} and $0$. I left \texttt{color} and \texttt{mycolor} non-deterministic, because~\cite{taubenfeld} says the initials do not matter.

Same-colour tickets use Lamport's pair order~\cite{lamport}, same as week~5: $(a,b)<(c,d)$ if $a<c$, or $a=c$ and $b<d$. Different colours: the smaller ticket is the one that is not the current \texttt{color} bit, so I kept the \texttt{if} on line~7. Line~3 is one atomic \texttt{Max} (the paper reads tickets one at a time; this is how week~5 writes the bakery maximum). No fairness on \texttt{ncs}, so a process can sit in its non-critical section forever (Dijkstra).

\paragraph{Mutex and Liveness.}
\begin{verbatim}
Mutex == [] \A i, k \in Procs :
           (i # k) => ~ ((pc[i] = "cs") /\ (pc[k] = "cs"))
Liveness == \A i \in Procs :
              []((pc[i] = "e1") => <>(pc[i] = "cs"))
\end{verbatim}
\texttt{Mutex} is the Peterson formula over three processes. \texttt{Liveness} uses \texttt{e1} because line~1 is when they show they want the CS. I did not write $\Box\diamond(\texttt{pc}[i]=\text{``cs''})$, because then they would also have to leave \texttt{ncs}.

TLC with \texttt{N <- 3} and \texttt{FairSpec} finished with no error ($992{,}752$ generated, $346{,}870$ distinct, diameter $103$). Both properties hold, it did not deadlock, and the \texttt{TypeOK} invariant confirms \texttt{number} never leaves \texttt{0..N} --- the bounded ticket space.

\paragraph{Fairness.}
\texttt{Mutex} is safety (week~5): if two processes are in \texttt{cs} together, a finite prefix already shows it. Fairness only talks about infinite runs, so it does not matter for Mutex. On the bare \texttt{Spec} (no $WF$/$SF$) TLC did not report a Mutex error.

\texttt{Liveness} is liveness, and without fairness it fails. \texttt{[][Next]\_vars} always lets the system stutter, so an enabled action can just never happen. TLC says \texttt{Liveness} is \emph{violated}. Last real state: $\texttt{pc}=\tuple{\texttt{e7},\texttt{e6},\texttt{e9}}$ (all three still waiting), then \texttt{Stuttering}.

I put fairness below the translation:
\begin{verbatim}
Active(self) == e1(self) \/ ... \/ x13(self)
FairSpec == Spec /\ \A self \in Procs : WF_vars(Active(self))
\end{verbatim}
\texttt{Active} is the process without \texttt{ncs}. Peterson puts $\texttt{SF\_vars}(P(0))\wedge\texttt{SF\_vars}(P(1))$ on the whole process, which is wrong here: \texttt{ncs} is always enabled, so fairness on $P(\textit{self})$ would force processes out of the remainder. $WF$ on \texttt{Active} still lets them sit in \texttt{ncs}, but once they do line~1 they have to keep going. $WF$ was enough; I did not need $SF$.

\paragraph{If one process breaks down.}
I put this in \texttt{bwbakeryfail.tla}, which \texttt{EXTENDS bwbakery}. A failed process goes to \texttt{pc[i] = "crashed"}. Nothing in Algorithm~2 is enabled from that label, so it takes no more steps. At most one crash, and no fairness on crashing. \texttt{MutexF} is the same as Mutex. \texttt{LivenessF} is Liveness but also allows \texttt{pc[i] = "crashed"}, so the one that dies is excused.

The answer depends on what you can still read from the crashed process.

\textbf{SpecZero} (same assumption as~\cite{lamport}): crash zeros \texttt{choosing[i]} and \texttt{number[i]}. TLC: no error ($2{,}547{,}516$ generated, $604{,}348$ distinct). Mutex and liveness still hold for the remaining two.

\textbf{SpecFreeze}: registers stay as they were. TLC did not report a \texttt{MutexF} error, but \texttt{LivenessF} is \emph{violated}: process~$2$ crashes with $\texttt{choosing}[2]=\texttt{TRUE}$, process~$1$ gets to line~6 with $j=2$ ($\texttt{pc}=\tuple{\texttt{e6},\texttt{crashed},\texttt{ncs}}$), then \texttt{Stuttering}. If someone dies in the doorway, \texttt{choosing} stays up and anyone at line~6 waits forever~\cite{lamport}.

So the remaining two only keep both properties if a failed process's registers eventually read as zero.

\begin{thebibliography}{9}
\bibitem{lamport} L.~Lamport. A New Solution of Dijkstra's Concurrent Programming Problem. \emph{Communications of the ACM}, 17(8), 1974.
\bibitem{taubenfeld} G.~Taubenfeld. The Black-White Bakery Algorithm and Related Bounded-Space, Adaptive, Local-Spinning and FIFO Algorithms. \emph{International Symposium on Distributed Computing}, 56--70, 2004.
\end{thebibliography}

\end{document}
